% Beweisversuch: ⟨δ_N, χ₃⟩ ∼ f(N) unter Hypothese H
\documentclass[11pt,a4paper]{article}

\usepackage[utf8]{inputenc}
\usepackage[T1]{fontenc}
\usepackage[ngerman]{babel}
\usepackage{amsmath,amssymb,amsthm}
\usepackage{booktabs}
\usepackage[a4paper,margin=2.5cm]{geometry}
\usepackage{microtype}
\usepackage{hyperref}

\title{Analytischer Beweisversuch:\\
$\langle\delta_N,\chi_3\rangle \sim f(N)$ unter Hypothese $H$}
\author{Thomas Hoffbauer}
\date{\today}

\newtheorem{theorem}{Theorem}
\newtheorem{proposition}[theorem]{Proposition}
\newtheorem{lemma}[theorem]{Lemma}
\newtheorem{hypothesis}{Hypothese}
\newtheorem{remark}{Bemerkung}
\newtheorem{conjecture}{Vermutung}

\newcommand{\R}{\mathcal{R}}
\newcommand{\E}{\mathbb{E}}
\newcommand{\Var}{\mathrm{Var}}

\begin{document}

\maketitle

\begin{abstract}
Wir formulieren eine präzise probabilistische Hypothese $H$ für die
Besetzung der sechs mod-$420$-Vierlingskanäle
$\R=\{11,101,191,221,311,401\}$,
leiten die explizite Skalierung
$f(N)=1/(6\sqrt{N})$ für die $\chi_3$-Projektion her und vergleichen sie
mit numerischen Daten bis $N_{\max}=10^9$.
Der kubische Charakter $\chi_3\bmod 7$ ist \emph{arithmetisch} auf
$\R$ mit dem $k{=}2$-DFT-Modus identisch ($\chi_3|_{\R}=\psi_2$);
unter $H$ ist ein nicht-null Koeffizient $a=\langle\delta_N,\chi_3\rangle$
eine \emph{statistische} Fluktuation der Größenordnung $O_p(1/\sqrt{N})$,
kein dynamischer Drift.
\end{abstract}

%=============================================================================
\section{Setup und Notation}
%=============================================================================

Sei $Q(p)=(p,p+2,p+6,p+8)$ ein Primzahlvierling mit Startprimzahl $p$.
Die sechs beobachteten Halbkanäle modulo $420$ sind
\[
  \R = \{11,101,191,221,311,401\}.
\]
Für $N$ gezählte Vierlingsereignisse mit Kanalzuordnung
$r_k\in\R$ sei $n_r(N)$ die Besetzungszahl und $K:=\sum_r n_r(N)=N$.

Der \emph{normalisierte Kanal-Bias} ist
\[
  \delta_N(r) \;:=\; \frac{n_r(N)}{N} - \frac{1}{6},
  \qquad r\in\R.
\]
Das Skalarprodukt auf $\mathbb{C}^{\R}$ ist
\[
  \langle f,g\rangle \;:=\; \frac{1}{6}\sum_{r\in\R} f_r\,\overline{g_r}.
\]
Der primitive kubische Dirichlet-Charakter $\chi_3\bmod 7$ ist durch
$\omega=e^{2\pi i/3}$ und
$\chi_3(4)=1$, $\chi_3(3)=\omega$, $\chi_3(2)=\omega^2$ gegeben.
Auf $\R$ gilt die exakte Identität (vgl.\ \texttt{eabc\_spectral\_decomposition.py}):
\[
  \chi_3|_{\R} = \psi_2,
  \qquad
  \psi_k(n) := e^{2\pi i kn/6},
\]
in der festen Reihenfolge $(11,101,191,221,311,401)$.

%=============================================================================
\section{Hypothese $H$ (formal)}
%=============================================================================

\begin{hypothesis}[$H$ --- Gleichverteilungsmodell auf $\R$]
\label{hyp:H}
Es gibt eine Konstante $C>0$, sodass die Vierlingsstarts mit
$p\to\infty$ unabhängig auftreten und die Kanalzuordnung bei $N$
Vierlingen durch ein \emph{Multinomialmodell}
\[
  (n_{11},n_{101},\ldots,n_{401}) \;\sim\;
  \mathrm{Mult}\!\left(N;\,\tfrac{1}{6},\ldots,\tfrac{1}{6}\right)
\]
beschrieben wird.
Äquivalent: die Folge der Kanal-Indizes
$(I_k)_{k=1}^{N}$ ist i.i.d.\ uniform auf $\{1,\ldots,6\}$,
und $n_r(N)=\#\{k: r_k=r\}$.
\end{hypothesis}

\subsection{Motivation: welches Modell ist natürlich?}

\paragraph{(i) Multinomial mit gleichen $p_i=1/6$.}
Dies ist das direkte diskrete Modell für $N$ unabhängige Kanalzuweisungen.
Es ist die kanonische Nullhypothese, wenn alle sechs Restklassen in $\R$
asymptotisch gleich häufig sind.

\paragraph{(ii) Poisson-Approximation.}
Für großes $N$ gilt
$\mathrm{Mult}(N;p)\approx (\mathrm{Poisson}(Np_1),\ldots,\mathrm{Poisson}(Np_6))$
mit gemeinsamer Summe $N$.
Die $\chi_3$-Projektion skaliert identisch, da
$\delta_N(r)=n_r/N-1/6$ nur von den relativen Anteilen abhängt.

\paragraph{(iii) Hardy--Littlewood-artige Dichte.}
Die HL-Vermutung für Vierlinge liefert
\[
  \#\{p\le x : Q(p)\text{ prim}\} \sim C_4 \int_2^x \frac{dt}{\log^4 t}.
\]
Die Singulärreihe $S(Q)$ hängt von lokalen Kongruenzbedingungen ab.
\emph{Wenn} die sechs Kanäle in $\R$ dieselbe lokale HL-Dichte tragen
(d.h.\ die Singulärreihen-Faktoren mod $420$ auf $\R$ übereinstimmen),
dann ist $H$ die asymptotische Konsequenz: bedingt auf ein Vierlingsereignis
ist der Start $p\bmod 420$ asymptotisch gleichverteilt auf $\R$.

\begin{remark}[Was $H$ \emph{nicht} behauptet]
$H$ postuliert keinen Chebyshev-artigen Drift
$B_{101,11}(N)/\sqrt{N}\to c\neq 0$,
keine Rang-1-Dynamik und keine kausale EABC-Interpretation.
Es ist ein \emph{statisches} Gleichverteilungsmodell für die Kanalzählung.
\end{remark}

\subsection{Arithmetische Struktur, die $\chi_3$-Korrelation erzwingt}

Die Korrelation mit $\chi_3$ ist \emph{algebraisch} erzwungen, nicht
probabilistisch:

\begin{lemma}[Mod-$7$-Zyklus und DFT-Identität]
\label{lem:cycle}
In der festen Reihenfolge der Kanäle gilt
$r_n\bmod 7 = (4,3,2,4,3,2)$, also
\[
  \chi_3(r_n) = (1,\omega,\omega^2,1,\omega,\omega^2) = \psi_2(n).
\]
Folglich ist der $\chi_3$-Unterraum \emph{identisch} mit dem
$k{=}2$-Fourier-Unterraum auf dem $6$-Ring:
\[
  \langle\delta_N,\chi_3\rangle = F_2(\delta_N)
  = \frac{1}{6}\sum_{n=0}^{5}\delta_n\,e^{-2\pi i\cdot 2n/6}.
\]
\end{lemma}

\begin{proof}
Direkte Rechnung; vgl.\ Lemma~2.1 in \texttt{eabc\_vierlinge\_kanalstruktur.tex}.
\end{proof}

\emph{Konsequenz:} Jeder Bias $\delta_N$ besitzt automatisch eine
$\chi_3$-Komponente; die Frage ist nur deren \emph{Größe}, nicht deren
\emph{Existenz}. Die arithmetische Struktur bestimmt die
\emph{Richtung} ($\chi_3|_{\R}$), das Modell $H$ bestimmt die
\emph{Skalierung} von $|\langle\delta_N,\chi_3\rangle|$.

Zusätzlich gilt auf allen $r\in\R$:
$r\equiv 1\pmod{5}$, daher $\chi_5|_{\R}\equiv 1$ (trivial);
$r\equiv 2\pmod{3}$, daher ist der quadratische Charakter mod $3$
ebenfalls trivial auf $\R$.

%=============================================================================
\section{Proposition und Beweisversuch}
%=============================================================================

\begin{proposition}[CLT-Skalierung für $\langle\delta_N,\chi_3\rangle$ unter $H$]
\label{prop:clt}
Unter Hypothese~\ref{hyp:H} gilt für $N\to\infty$:
\begin{enumerate}
  \item $\E\bigl[\langle\delta_N,\chi_3\rangle\bigr] = 0$;
  \item $\Var\bigl(\langle\delta_N,\chi_3\rangle\bigr) = \dfrac{1}{36N}$;
  \item $\sqrt{N}\,\langle\delta_N,\chi_3\rangle
        \;\xRightarrow{d}\;
        \mathcal{N}_{\mathbb{C}}(0,\tfrac{1}{36})$
        (komplexe Normalverteilung);
  \item insbesondere
        $|\langle\delta_N,\chi_3\rangle| = O_p\!\bigl(1/\sqrt{N}\bigr)$,
        und für jedes $\varepsilon>0$:
        $|\langle\delta_N,\chi_3\rangle| = O_p\!\bigl(\sqrt{\log\log N/N}\bigr)$
        (LIL-artige obere Schranke für Partialsummen).
\end{enumerate}
\end{proposition}

\begin{proof}[Beweisversuch]
\emph{Schritt 1: Darstellung als gewichtete Zählung.}
Da $\sum_{r\in\R}\chi_3(r)=2(1+\omega+\omega^2)=0$, kürzt der Mittelwert:
\[
  \langle\delta_N,\chi_3\rangle
  = \frac{1}{6}\sum_r \Bigl(\frac{n_r}{N}-\frac{1}{6}\Bigr)\overline{\chi_3(r)}
  = \frac{1}{6N}\sum_{r\in\R} n_r\,\overline{\chi_3(r)}.
\]
Setze $S_N := \sum_{k=1}^{N}\overline{\chi_3(r_{I_k})}$.
Dann $\langle\delta_N,\chi_3\rangle = S_N/(6N)$.

\emph{Schritt 2: Erwartungswert.}
Unter $H$ ist $I_k$ uniform, also
$\E[\overline{\chi_3(r_{I_k})}]
=\frac{1}{6}\sum_r \overline{\chi_3(r)}=0$,
daher $\E[S_N]=0$ und $\E[\langle\delta_N,\chi_3\rangle]=0$.

\emph{Schritt 3: Varianz.}
Die $Y_k:=\overline{\chi_3(r_{I_k})}$ sind i.i.d.\ mit
$\E[Y_k]=0$ und $\E[|Y_k|^2]=\frac{1}{6}\sum_r|\chi_3(r)|^2
=\frac{1}{6}\cdot 6 = 1$.
Daher
\[
  \Var(S_N) = N\,\Var(Y_1) = N,
  \qquad
  \Var\!\left(\frac{S_N}{6N}\right) = \frac{N}{36N^2} = \frac{1}{36N}.
\]

\emph{Schritt 4: Zentraler Grenzwertsatz.}
Nach dem multidimensionalen ZGWS gilt
$S_N/\sqrt{N} \xRightarrow{d} \mathcal{N}_{\mathbb{C}}(0,1)$,
also
$\sqrt{N}\,\langle\delta_N,\chi_3\rangle
= S_N/(6\sqrt{N}) \xRightarrow{d} \mathcal{N}_{\mathbb{C}}(0,1/36)$.

\emph{Schritt 5: Größenordnung.}
Aus (3) folgt
$|\langle\delta_N,\chi_3\rangle| = O_p(1/\sqrt{N})$.
Die LIL-Aussage folgt aus dem Gesetz des iterierten Logarithmus
für i.i.d.\ komplexwertige Partialsummen $S_N$
(Kolmogorov--LIL): für a.s.\ alle $N$ groß genug,
$|S_N| \le (2N\log\log N)^{1/2}(1+o(1))$,
also $|\langle\delta_N,\chi_3\rangle|
\le (2\log\log N/(36N))^{1/2}(1+o(1))$.
\end{proof}

\subsection{Explizite Funktion $f(N)$}

\begin{corollary}[Explizite Skalierung]
\label{cor:fN}
Unter $H$ ist die natürliche Skalierungsfunktion
\[
  \boxed{
  f(N) = \frac{1}{6\sqrt{N}}
  }
\]
im Sinne von
\[
  \langle\delta_N,\chi_3\rangle
  \;\sim\;
  \mathcal{N}_{\mathbb{C}}\!\left(0,\,\frac{1}{36N}\right),
\]
d.h.\ typische Größe $|\langle\delta_N,\chi_3\rangle|\asymp f(N)$.

Präziser für den Betrag (komplexe Normalverteilung):
\[
  \E\bigl[|\langle\delta_N,\chi_3\rangle|\bigr]
  \;\approx\;
  \sqrt{\frac{\pi}{72N}}
  \;=\;
  \frac{\sqrt{\pi/2}}{6\sqrt{N}}
  \;\approx\;
  0{,}886\cdot f(N).
\]
Der $z$-Score $|\langle\delta_N,\chi_3\rangle|/f(N)$ ist unter $H$
asymptotisch $\mathrm{Rayleigh}(1)$-verteilt (Median $\approx 1{,}177$).
\end{corollary}

%=============================================================================
\section{Numerische Validierung}
%=============================================================================

\subsection{$\chi_3$-Projektion bei drei Skalen}

Tabelle~\ref{tab:chi3} stammt aus \texttt{eabc\_analytic\_proof\_check.py}
(Multinomial-Null: $10\,000$ Trials).

\begin{table}[h]
\centering
\caption{$\chi_3$-Skalierung vs.\ Theorie $f(K)=1/(6\sqrt{K})$.}
\label{tab:chi3}
\begin{tabular}{rrrrrrr}
\toprule
$N_{\max}$ & $K$ & $|\langle\delta,\chi_3\rangle|$ & $f(K)$ & $z$ &
$|\langle\delta,\chi_3\rangle|\sqrt{K}$ & $p(\ge)$ \\
\midrule
$10^8$    & 4767  & 0.002749 & 0.002414 & 1.14 & 0.190 & 26.8\% \\
$3\cdot10^8$ & 10972 & 0.001568 & 0.001591 & 0.99 & 0.164 & 37.7\% \\
$10^9$    & 28387 & 0.001043 & 0.000989 & 1.05 & 0.176 & 33.1\% \\
\bottomrule
\end{tabular}
\end{table}

\paragraph{Urteil.}
\begin{itemize}
  \item $|\langle\delta,\chi_3\rangle|$ fällt exakt wie $1/\sqrt{K}$
        (Produkt $|\langle\delta,\chi_3\rangle|\sqrt{K}\approx 0.18\pm 0.01$).
  \item Alle $z$-Scores liegen nahe $1$ (erwarteter Median $\approx 1.18$).
  \item Null-$p$-Werte ($27$--$38\%$) zeigen: kein signifikanter
        systematischer $\chi_3$-Bias jenseits der Multinomial-Fluktuation.
\end{itemize}

\subsection{Vergleich mit $B_{101,11}/\sqrt{K}$}

Aus \texttt{eabc\_scaling\_test.py}:

\begin{table}[h]
\centering
\caption{Chebyshev-artiger Paar-Bias $B_{101,11}=N_{101}-N_{11}$.}
\label{tab:B}
\begin{tabular}{rrrr}
\toprule
$N_{\max}$ & $K$ & $B_{101,11}$ & $B/\sqrt{K}$ \\
\midrule
$10^8$    & 4767  & $+66$ & $+0.956$ \\
$3\cdot10^8$ & 10972 & $+51$ & $+0.487$ \\
$10^9$    & 28387 & $+77$ & $+0.457$ \\
\bottomrule
\end{tabular}
\end{table}

\paragraph{Warum fällt $B/\sqrt{K}$?}
$B_{101,11}$ ist ein \emph{Pfadprozess} (Random Walk):
jeder Vierling trägt $+1$ (Kanal $101$), $-1$ (Kanal $11$) oder $0$.
Unter $H$:
$\Var(B_K)=K/3$, also $B_K/\sqrt{K}$ hat Varianz $1/3$
(stationär verteilt, kein Limes).
Der beobachtete Abfall $0.96\to 0.46$ ist \emph{ein einzelner Pfad}:
die Null-$p$-Werte steigen von $9.8\%$ auf $43\%$,
konsistent mit Random-Walk-Fluktuation, nicht mit Kollaps auf $0$.

\emph{Wichtig:} $B/\sqrt{K}$ und $|\langle\delta,\chi_3\rangle|$ messen
verschiedene Projektionen:
$B$ projiziert auf $\chi_4$-Richtung (Kanalpaar $11$ vs.\ $101$),
$\langle\delta,\chi_3\rangle$ auf den kubischen $3$-Wellen-Unterraum.

%=============================================================================
\section{Mechanismus für $a\neq 0$}
%=============================================================================

Die Zerlegung
\[
  \delta_N = a\,\hat{e}_{\chi_3} + \bar{a}\,\hat{e}_{\bar{\chi}_3}
           + b\,\hat{e}_{\chi_4} + c\,\hat{e}_{\chi_4\chi_3} + \varepsilon
\]
hat Koeffizient $a=\langle\delta_N,\hat{e}_{\chi_3}\rangle$.
Bei $N=4767$: $|a|\approx 0.00275$, Varianzanteil $\chi_3+\bar{\chi}_3$: $59.7\%$.

\subsection{Warum $a\neq 0$, obwohl alle Kanäle ``gleich'' sind?}

\paragraph{Antwort unter $H$:} $a\neq 0$ ist \emph{typisch}.
Unter $\E[a]=0$ und $\Var(a)\sim 1/N$ ist
$\mathbb{P}(a\neq 0)\to 1$;
die Frage ist nur, ob $|a|$ größer als $O(1/\sqrt{N})$ ist.
Bei $K=4767$: $z=1.14$, also \emph{kein} überraschender Befund.

\paragraph{Arithmetische vs.\ statistische Ebene.}
\begin{itemize}
  \item \textbf{Arithmetik} erzwingt: die Projektions\emph{richtung} $\chi_3$
        ist fest (Lemma~\ref{lem:cycle}). Jede Abweichung von Gleichverteilung
        hat automatisch eine $\chi_3$-Komponente.
  \item \textbf{Statistik} unter $H$: die \emph{Größe} $|a|$ ist
        $O_p(1/\sqrt{N})$, nicht konstant.
  \item Die hohe Varianzquote $59.7\%$ bei $K=4767$ bedeutet nicht
        ``$59\%$ des Bias ist $\chi_3$-Ursache'', sondern:
        \emph{von der beobachteten Abweichung $\delta_N$}
        entfällt $59\%$ auf den $\chi_3$-Unterraum --- bei
        zufälligem $\delta_N$ wäre dies unter $H$ ebenfalls
        $\approx 50$--$60\%$ (da $\chi_3$ zwei reelle Dimensionen
        von sechs trägt und $|{\chi_3(r)}|=1$).
\end{itemize}

\subsection{Bedingte Verteilung $p\bmod 7\,|\,\text{Vierling}$}

Die Kanalzuordnung hängt von $p\bmod 420$ ab, nicht nur von $p\bmod 7$.
Auf $\R$ gibt es nur drei $\chi_3$-Werte mit Multiplizität $2$:
\[
  \#\{r\in\R: r\bmod 7 = 4\} = 2,\quad
  \#\{r\in\R: r\bmod 7 = 3\} = 2,\quad
  \#\{r\in\R: r\bmod 7 = 2\} = 2.
\]
Unter exakter Gleichverteilung auf $\R$ ist die bedingte Verteilung
$p\bmod 7\,|\,Q(p)$ gleich auf $\{2,3,4\}$ (je $1/3$).
Ein $\chi_3$-Bias entstünde nur, wenn die \emph{tatsächliche}
Vierlingsdichte auf den drei mod-$7$-Klassen \emph{ungleich} wäre.

\emph{Beweisversuch-Lücke:} Die Gleichverteilung der HL-Dichte auf
$\R$ (nicht nur auf $\{2,3,4\}\bmod 7$) ist nicht bewiesen.
Numerisch weichen die Kanalzählungen bei $K=4767$ nur um $\pm 1\%$
vom Mittel ab --- konsistent mit $O(1/\sqrt{K})$-Fluktuation,
nicht mit strukturellem mod-$7$-Bias.

%=============================================================================
\section{Was ist bewiesen vs.\ offen}
%=============================================================================

\subsection{Bewiesen (exakt oder numerisch verifiziert)}

\begin{enumerate}
  \item $\chi_3|_{\R}=\psi_2$ (algebraische Identität).
  \item $F_2=\langle\delta_N,\chi_3\rangle$ (Theorem~A).
  \item Vollständige Charakterzerlegung von $\delta_N$ bei $K=4767$
        (Residuum $<10^{-36}$).
  \item Unter $H$: $\Var(\langle\delta_N,\chi_3\rangle)=1/(36N)$
        (Proposition~\ref{prop:clt}, Beweisversuch vollständig
        \emph{conditional on $H$}).
  \item Numerische Skalierung $|\langle\delta,\chi_3\rangle|\sim 1/\sqrt{K}$
        über drei Dekaden (Tabelle~\ref{tab:chi3}).
  \item $B/\sqrt{K}$ zeigt keinen stabilen Drift (Tabelle~\ref{tab:B}).
\end{enumerate}

\subsection{Offene Lücken (für vollständigen Beweis)}

\begin{enumerate}
  \item \textbf{HL-Äquidistribution auf $\R$:}
        Beweis, dass die Singulärreihe für Vierlinge auf den
        sechs Kanälen in $\R$ asymptotisch gleich ist.
  \item \textbf{Unabhängigkeit:}
        $H$ postuliert i.i.d.\ Kanalzuweisungen; echte Vierlinge
        haben langreichweitige arithmetische Korrelationen
        (Gemeinsamkeit von Primzahlpaaren). Für $N\sim 10^4$
        ist dies vermutlich vernachlässigbar, aber nicht bewiesen.
  \item \textbf{Uniformitätsabschätzung:}
        Explizite Fehlerterm-Schranke
        $|\langle\delta_N,\chi_3\rangle - O_p(1/\sqrt{N})|
        = O(N^{-1/2-\varepsilon})$ aus Primzahlverteilung in
        arithmetischen Progressionen.
  \item \textbf{Signifikanz:}
        Ein $z\approx 1$ ist nicht signifikant; für Behauptung
        ``struktureller $\chi_3$-Bias'' müsste $z\to\infty$ mit $N$
        nachgewiesen werden --- dies geschieht \emph{nicht}
        (Tabelle~\ref{tab:chi3}).
\end{enumerate}

%=============================================================================
\section{Epistemologische Abgrenzung}
%=============================================================================

\subsection{Diagonalisierung vs.\ Ursache}

\begin{itemize}
  \item \textbf{Diagonalisierung (bewiesen):}
        $\delta_N$ lässt sich in Charaktere zerlegen;
        $\chi_3$ ist ein \emph{Koordinatensystem}, kein \emph{Mechanismus}.
        Die Identität $\chi_3|_{\R}=\psi_2$ ist reine Fourieranalysis
        auf der endlichen Menge $\R$.

  \item \textbf{Ursache (offen):}
        \emph{Warum} weichen die Kanalzählungen von $N/6$ ab?
        Unter $H$: Zufallsfluktuation.
        Gegen $H$: HL-Singularreihen-Unterschiede, Korrelationen
        zwischen benachbarten Vierlingen, oder dynamische Prozesse
        (EABC, Chebyshev-Bias).

  \item \textbf{Negativer Skalierungstest:}
        $B/\sqrt{K}$ kollabiert nicht auf $0$, aber schwankt
        wie ein Random Walk --- kein Beweis für persistenten Drift.
        Die $\chi_3$-Projektion skaliert korrekt als $1/\sqrt{K}$,
        was \emph{gegen} einen anwachsenden strukturellen Bias spricht.
\end{itemize}

\subsection{Zusammenfassung der Hypothese und $f(N)$}

\begin{center}
\fbox{\parbox{0.92\textwidth}{
\textbf{Hypothese $H$:}
Vierlingskanäle sind i.i.d.\ uniform auf $\R$ (Multinomial $1/6$).

\medskip
\textbf{Explizite $f(N)$:}
$\displaystyle f(N)=\frac{1}{6\sqrt{N}}$,
$\langle\delta_N,\chi_3\rangle\sim\mathcal{N}_{\mathbb{C}}(0,1/(36N))$.

\medskip
\textbf{Daten:} Konsistent bei $K=4767,10972,28387$;
$z\approx 1$, $|\langle\delta,\chi_3\rangle|\sqrt{K}\approx 0.18$.

\medskip
\textbf{Status:} Beweisversuch \emph{conditional on $H$} vollständig;
$H$ selbst (HL-Äquidistribution) offen.
}}
\end{center}

%=============================================================================
\appendix
\section{Reproduzierbarkeit}

\begin{itemize}
  \item \texttt{eabc\_analytic\_proof\_check.py} --- $\chi_3$-Skalierung
  \item \texttt{eabc\_scaling\_test.py} --- $B_{101,11}/\sqrt{K}$
  \item \texttt{eabc\_spectral\_decomposition.py} --- Zerlegung bei $K=4767$
\end{itemize}

\end{document}
